\documentclass[zihao=-4,a4paper]{ctexart}
\usepackage[margin=2.4cm]{geometry}
\usepackage{amsmath,amssymb,booktabs,graphicx,float,array,enumitem,xcolor}
\usepackage[most]{tcolorbox}
\usepackage[hidelinks]{hyperref}
\linespread{1.35}
\setlist{itemsep=2pt,topsep=4pt}
\ctexset{section={format=\Large\bfseries\color{blue!55!black}},subsection={format=\large\bfseries}}
\newtcolorbox{keybox}[1][本次要点]{colback=blue!3,colframe=blue!45!black,fonttitle=\bfseries,title=#1,boxrule=0.6pt,arc=1mm}
\newtcolorbox{notebox}{colback=gray!6,colframe=gray!60,boxrule=0.4pt,arc=1mm}
\newcommand{\fig}[3][0.95]{\begin{figure}[H]\centering\includegraphics[width=#1\linewidth]{#2}\caption{#3}\end{figure}}
\newcommand{\M}{M^{(K)}}

\title{\textbf{第 2 次组会汇报}\\[4pt]\Large $\M$ 比现有做法多看到了什么}
\author{对应工作：103—106、108、109、117 号实验}
\date{}

\begin{document}
\maketitle
\thispagestyle{empty}

\begin{keybox}
\begin{enumerate}[leftmargin=*]
  \item 先把“攻击者能推出多少”的口径定死：\textbf{三种攻击器取最大 + 单调闭包}。只用一个 DNN 会低估。
  \item 完整风险表显示，只看单字段时，$\tau=0.7$ 下关键条目仅占 2\%—4\%；允许 2 个背景字段后升到 50\%—68\%。
  \item 与相关系数、SAGE、LOCO、推断图等现有打分相比，$M^{(2)}$ 不会漏掉协同、冗余和间接三类结构；决策上，单字段规则\textbf{只找到 13\%—17\% 的关键字段}。
  \item 落到发布审查：单字段规则会放行 19\% / 56\% 的危险集合，按 $M$ 做预算则\textbf{零危险放行}。
  \item 两条诚实的边界：单字段推断能力做粗排并不差；$\M$ 是最大值统计量，小样本和弱攻击者都会让结论出偏。
\end{enumerate}
\end{keybox}

\section{正式口径：攻击者能推出多少（103 号）}

真值 $V(S)$ 由三种攻击器（单目标 DNN、多目标 DNN、梯度提升树）逐集合专用训练得到，在验证集上选最强的一个，再做单调闭包。

\fig{figures/图1.png}{103 号：各攻击器被选中的占比（左），以及只用一个 DNN 时的低估幅度（右）\protect\newline{\footnotesize 原图：\texttt{DNN\_Aggresvation103/figures/fig3\_attacker\_share.png}}}

\begin{itemize}[leftmargin=*]
  \item 集合规模为 3 时，梯度提升树被选中 55\%（PJM）/ 43\%（CAISO），多目标 DNN 只占约 10\%。
  \item 只用一个 DNN 会让 $V$ 低估 0.005—0.027，且集合越大低估越多。\textbf{100 号及之前的数值因此偏低}，后面一律使用本口径。
\end{itemize}

\section{完整风险表：风险主要来自背景（104 号）}

PJM、CAISO 各 41 个字段 $\times$ 12 个目标，$K=0,1,2$ 全部精确计算。

\fig{figures/图2.png}{104 号：正式真值下的推断升级曲线\protect\newline{\footnotesize 原图：\texttt{DNN\_Aggresvation104/figures/fig3\_escalation\_official.png}}}

\begin{table}[H]\centering
\caption{$\tau=0.7$ 时的关键（字段，目标）条目（共 492 条）}
\begin{tabular}{lccc}
\toprule
& $K=0$（只看单字段） & $K=1$ & $K=2$ \\
\midrule
PJM & 11（2.2\%） & 115 & \textbf{337（68\%）} \\
CAISO & 18（3.7\%） & 137 & \textbf{246（50\%）} \\
\bottomrule
\end{tabular}
\end{table}

每个关键条目都带有见证背景，例如：

\begin{table}[H]\centering
\caption{见证背景示例（精确真值）}
\footnotesize\setlength{\tabcolsep}{4pt}
\begin{tabular}{llcccl}
\toprule
目标 & 字段 & $M^{(0)}$ & $M^{(1)}$ & 见证背景 & 机制 \\
\midrule
PJM 总发电 & 风电占比 & 0.162 & 0.871 & 风电出力 & 总发电 = 出力 $\div$ 占比 \\
PJM 总发电 & 光伏出力 & 0.033 & 0.643 & 光伏占比 & 同上 \\
CAISO 实时 SP15 阻塞价 & 实时 NP15 电价 & 0.040 & 0.406 & 实时 SP15 电价 & 阻塞 $\approx$ 节点价差 \\
CAISO 实际负荷 & 南区负荷计划 & 0.319 & 0.635 & SP15 光伏预测 & 净负荷关系 \\
\bottomrule
\end{tabular}
\end{table}

\section{与现有字段打分方法对比（105、117 号）}

\subsection{机理：三类方法各有盲区}

构造一个结构已知的小例子 $Y=a\cdot A\cdot B+b\cdot P+\text{噪声}$：$A$、$B$ 协同（单独无用），$C_1$、$C_2$ 是 $P$ 的两个带噪副本（冗余），$D$ 是间接字段（$A+B$ 为 0.54，加上 $D$ 后为 0.73，越过 0.7）。

\fig{figures/图3.png}{117 号：各方法在协同、冗余、间接字段上的分数（63 个集合全部重训）\protect\newline{\footnotesize 原图：\texttt{DNN\_Aggresvation117/figures/fig1\_toy\_mechanism.png}}}

\begin{itemize}[leftmargin=*]
  \item 相关系数、单字段能力、推断图（成对边）\textbf{看不见协同}：$A$、$B$ 都是 0。
  \item LOCO、置换重要性、SAGE 在全字段模型上计算，\textbf{冗余字段互为替身}（$C_1$ 的 LOCO 为 0.01），\textbf{间接字段 $D$ 被掩盖为 0}。
  \item $M^{(2)}$ 在三类结构上都给出正确的量级。
\end{itemize}

\subsection{检测关键字段：单字段推断能力是强基线}

\fig{figures/图4.png}{105 号：$\tau=0.7$、$K=2$ 下检测关键字段的 PR-AUC\protect\newline{\footnotesize 原图：\texttt{DNN\_Aggresvation105/figures/fig1\_baseline\_pr.png}}}

\begin{notebox}
\textbf{对本文不利的结果，如实报告}：只做“排序”时，单字段重训能力 $M^{(0)}$ 在 PJM 上的 PR-AUC 为 0.877，与 $M^{(2)}$ 持平；只有在 CAISO 上 $M$ 明显更好（$0.764\to0.829$）。原因是单字段能力为 0.6 的字段，加一个背景就容易越过 0.7。\\
所以论文不能说“单字段分析完全看不到组合风险”。准确的说法是：单字段能力能粗排出“可能危险的字段”，但给不出\textbf{危险有多大、在什么背景下危险、是否已越过阈值}。模型归因类方法（SAGE、置换重要性）最差。
\end{notebox}

\subsection{决策层面：差距才真正拉开}

117 号在 10 个目标（第 3 次汇报会介绍的五组数据）上比较“按规则定出的关键字段”与“扣留字段后还剩多少危险组合”。

\fig{figures/图5.png}{117 号：关键字段召回（左）与扣留后仍未打破的危险组合（右）。估计器为当时版本，最终版数字见第 5 次\protect\newline{\footnotesize 原图：\texttt{DNN\_Aggresvation117/figures/fig2\_decision\_disposal.png}}}

\begin{itemize}[leftmargin=*]
  \item 现行的单字段定级规则只召回 \textbf{13\%（$\tau=0.5$）/ 17\%（$\tau=0.7$）}的关键字段。
  \item 按单字段定级扣留，扣得少，但 \textbf{98.2\% / 95.5\% 的危险组合仍未被打破}；按 $\hat M^{(2)}$ 自适应扣留，只多扣约 1 个字段，残余危险组合降到 2\% 左右。
\end{itemize}

\section{落到发布审查（106 号）}

随机抽 2,000 个规模为 3 的待发布集合，$\tau=0.7$，用正式真值判断它们实际是否安全。

\fig{figures/图6.png}{106 号：各审查规则的危险放行率（左）与误拒率（右）\protect\newline{\footnotesize 原图：\texttt{DNN\_Aggresvation106/figures/fig1\_release\_audit.png}}}

\begin{table}[H]\centering
\caption{发布审查结果（PJM / CAISO）}
\begin{tabular}{lcc}
\toprule
审查规则 & 危险放行 & 误拒 \\
\midrule
单字段能力 $\le\tau$ 即放行 & \textbf{19.1\% / 56.4\%} & 0 / 0 \\
通用模型直接估计 $\hat V(A)\le\tau$ & 6.6\% / 2.6\% & 0 / 0 \\
逐阶预算 $U^{(2)}\le\tau$ & \textbf{0 / 0} & 50\% / 33\% \\
精确 $M$ 预算 $\sum M^{(2)}\le\tau$ & 0 / 0 & 86\% / 74\% \\
\bottomrule
\end{tabular}
\end{table}

\begin{itemize}[leftmargin=*]
  \item 这是最直接的安全论据：按 $M$ 做预算有理论保证（第 1 次的预算定理），\textbf{零危险放行}。
  \item 逐阶预算 $U^{(2)}$ 是实用的折中：仍然零危险放行，误拒从 86\% / 74\% 降到 50\% / 33\%。
  \item 由此得到两级流程：先用逐阶预算快速放行，没通过的再用直接估计加重训认证复核。
\end{itemize}

\section{两条边界}

\subsection{用估计值判定时要给区间（108 号）}

用通用模型估计的 $\hat M$ 做点判定，PJM $K=2$、$\tau=0.7$ 时关键字段召回只有 0.754。改为区间 $[L,U]$：下界 $L$ 取前 3 个候选见证背景的重训值，上界 $U$ 用留一字段校准的分位数。然后做保守判定。

\fig{figures/图7.png}{108 号：点估计与区间保守判定的对比\protect\newline{\footnotesize 原图：\texttt{DNN\_Aggresvation108/figures/fig2\_critical\_interval.png}}}

召回从 0.754 升到 0.994，代价是精确率从约 0.99 降到 0.87（误报交给重训认证复核）。$\tau=0.9$ 时精确率明显下降，需要单独说明。

\subsection{攻击者能力的假设很关键（109 号）}

让攻击者只拿到 10\%—50\% 的带标签辅助数据：

\begin{itemize}[leftmargin=*]
  \item 推断能力下降有限：只用 10\% 数据，三字段组合的平均 $V$ 只下降 10\%—19\%。\textbf{数据少并不是有效的防护。}
  \item 反直觉的是，小样本下 $M^{(2)}$ 反而比全量高 5\%—12\%。原因是 $\M$ 是\textbf{最大值统计量}：候选背景越多（$K=2$ 时有 821 个），越容易挑中正噪声（winner's curse）。
  \item 审计方如果按弱攻击者评估，会单向漏判：按 10\% 数据评估并放行时，PJM 在 $\tau=0.5$ 下有 \textbf{25.6\%} 的危险集合被放行。
\end{itemize}

\fig{figures/图8.png}{109 号：选择增益随候选背景数增大（左）；$K=2$ 的偏差分解（中）；小样本挑出的见证背景大多是噪声（右）\protect\newline{\footnotesize 原图：\texttt{DNN\_Aggresvation109/figures/fig4\_winners\_curse.png}}}

结论：审计必须假设最强的攻击者（攻击器族取最大）；报告 $M$ 数值时优先用重训认证过的下界，而不是估计值取最大后的结果。

\section{小结与下一次}

\begin{keybox}[小结]
\begin{itemize}[leftmargin=*]
  \item $\M$ 的价值主要不在排序，而在\textbf{决策}：它能说出危险有多大、在什么背景下危险、是否越线，并带有零危险放行的保证。
  \item 单字段规则只找到 13\%—17\% 的关键字段，并放行 19\% / 56\% 的危险集合。
  \item 使用时要注意：给区间而不是点估计；按最强攻击者审计；数值要考虑取最大值带来的偏高。
\end{itemize}
\end{keybox}

\textbf{下一次}：前面都是 PJM、CAISO 的系统级字段。下一次换到真实电力场景（标准算例、澳洲 NEM 的真实机组、按披露时点重新定义的 PJM / CAISO），看组合风险是否依然存在、有多大。

\end{document}
